\chapter{第8章：Eichler--Shimura 习题}
\difficultykey

\section{基础题 \easy}

\begin{problem}{小素数处的 Hecke 系数}
\easy
对导子 $11$ 的椭圆曲线
\[
E_{11}: y^2+y=x^3-x^2-10x-20
\]
利用点数公式 $a_p=p+1-\#E_{11}(\mathbb{F}_p)$（对好素数 $p$）计算 $a_2,a_3,a_5,a_7$。
\end{problem}
\begin{solution}
由第 7 章计算得
\[
\#E_{11}(\mathbb{F}_2)=5,
\quad
\#E_{11}(\mathbb{F}_3)=5,
\quad
\#E_{11}(\mathbb{F}_5)=5,
\quad
\#E_{11}(\mathbb{F}_7)=10.
\]
这些素数都不整除导体 $11$，因此都是好素数。于是
\[
a_2=2+1-5=-2,
\qquad
a_3=3+1-5=-1,
\]
\[
a_5=5+1-5=1,
\qquad
a_7=7+1-10=-2.
\]
所以
\[
(a_2,a_3,a_5,a_7)=(-2,-1,1,-2).
\]
\end{solution}

\begin{problem}{Euler 因子}
\easy
写出 $L(E_{11},s)$ 在素数 $p=2,3,5,7,11$ 处的局部 Euler 因子。
\end{problem}
\begin{solution}
对好素数 $p\nmid 11$，局部因子为
\[
L_p(E_{11},s)=\bigl(1-a_pp^{-s}+p^{1-2s}\bigr)^{-1}.
\]
因此
\[
L_2(E_{11},s)=\bigl(1+2\cdot 2^{-s}+2\cdot 2^{-2s}\bigr)^{-1},
\]
\[
L_3(E_{11},s)=\bigl(1+3^{-s}+3\cdot 3^{-2s}\bigr)^{-1},
\]
\[
L_5(E_{11},s)=\bigl(1-5^{-s}+5\cdot 5^{-2s}\bigr)^{-1},
\]
\[
L_7(E_{11},s)=\bigl(1+2\cdot 7^{-s}+7\cdot 7^{-2s}\bigr)^{-1}.
\]
在 $p=11$ 处，$E_{11}$ 是分裂乘法约化，因此局部因子为
\[
L_{11}(E_{11},s)=\bigl(1-11^{-s}\bigr)^{-1}.
\]
\end{solution}

\begin{problem}{Frobenius 特征多项式}
\easy
设 $p\nmid N$。从 Eichler--Shimura 关系
\[
T_p\equiv \mathrm{Frob}_p+p\,\mathrm{Frob}_p^{-1}
\]
说明：若某个 Hecke 本征向量上 $T_p$ 的本征值是 $a_p$，则 Frobenius 的特征多项式为
\[
X^2-a_pX+p.
\]
\end{problem}
\begin{solution}
设 Frobenius 在该二维空间上的两个特征值为 $\alpha,\beta$。因为它们互为倒数乘以 $p$，所以
\[
\alpha\beta=p.
\]
而 Eichler--Shimura 关系说明 $T_p$ 的作用等于 $\alpha+\beta$，于是
\[
\alpha+\beta=a_p.
\]
因此 Frobenius 的特征多项式就是
\[
(X-\alpha)(X-\beta)=X^2-(\alpha+\beta)X+\alpha\beta=X^2-a_pX+p.
\]
\end{solution}

\begin{problem}{二次扩张上的点数}
\easy
利用上一题中的 Frobenius 特征值，证明对好素数 $p$ 有
\[
\#E(\mathbb{F}_{p^2})=p^2+1-(a_p^2-2p),
\]
并计算 $\#E_{11}(\mathbb{F}_4)$ 与 $\#E_{11}(\mathbb{F}_{25})$。
\end{problem}
\begin{solution}
记 Frobenius 特征值为 $\alpha,\beta$，则
\[
\alpha+\beta=a_p,
\qquad
\alpha\beta=p.
\]
在 $\mathbb{F}_{p^2}$ 上，迹变成
\[
\alpha^2+\beta^2=(\alpha+\beta)^2-2\alpha\beta=a_p^2-2p.
\]
故
\[
\#E(\mathbb{F}_{p^2})=p^2+1-(\alpha^2+\beta^2)=p^2+1-(a_p^2-2p).
\]
对 $E_{11}$，
\[
a_2=-2 \Rightarrow \#E_{11}(\mathbb{F}_4)=4+1-\bigl((-2)^2-4\bigr)=5,
\]
\[
a_5=1 \Rightarrow \#E_{11}(\mathbb{F}_{25})=25+1-(1-10)=35.
\]
\end{solution}

\begin{problem}{系数递推与乘法性}
\easy
利用
\[
a_{mn}=a_ma_n\quad ((m,n)=1),
\qquad
a_{p^2}=a_p^2-p\quad (p\nmid 11)
\]
计算 $a_4,a_6,a_9,a_{10}$。
\end{problem}
\begin{solution}
由上一节的 $a_2=-2,a_3=-1,a_5=1$，得到
\[
a_4=a_2^2-2=4-2=2,
\qquad
a_9=a_3^2-3=1-3=-2.
\]
又因为 $2,3,5$ 两两互素，
\[
a_6=a_2a_3=(-2)(-1)=2,
\qquad
a_{10}=a_2a_5=(-2)(1)=-2.
\]
所以
\[
a_4=2,\quad a_6=2,\quad a_9=-2,\quad a_{10}=-2.
\]
\end{solution}

\begin{problem}{前几项 Dirichlet 级数}
\easy
写出 $L(E_{11},s)=\sum_{n\ge 1}a_nn^{-s}$ 的前十项。
\end{problem}
\begin{solution}
由 $a_1=1$ 以及前几题得到
\[
a_2=-2,\ a_3=-1,\ a_4=2,\ a_5=1,\ a_6=2,\ a_7=-2,\ a_8=0,\ a_9=-2,\ a_{10}=-2.
\]
在坏素数 $11$ 处，分裂乘法约化给出 $a_{11}=1$。因此
\[
L(E_{11},s)=1-\frac{2}{2^s}-\frac{1}{3^s}+\frac{2}{4^s}+\frac{1}{5^s}+\frac{2}{6^s}-\frac{2}{7^s}
-\frac{2}{9^s}-\frac{2}{10^s}+\frac{1}{11^s}+\cdots.
\]
其中 $a_8=0$，所以没有 $8^{-s}$ 项。
\end{solution}

\begin{problem}{完成的 L 函数}
\easy
写出 $E_{11}$ 的完成 $L$-函数 $\Lambda(E_{11},s)$，并指出函数方程的中心点。
\end{problem}
\begin{solution}
由于导体 $N_{E_{11}}=11$，权重为 $2$，完成 $L$-函数是
\[
\Lambda(E_{11},s)=11^{s/2}(2\pi)^{-s}\Gamma(s)L(E_{11},s).
\]
Eichler--Shimura 和 Weil 定理说明它满足形如
\[
\Lambda(E_{11},s)=\varepsilon\,\Lambda(E_{11},2-s),
\qquad \varepsilon=\pm 1
\]
的函数方程。对权 $2$ 情形，对称中心是
\[
s=1.
\]
这也是 BSD 猜想关注的特殊点。
\end{solution}

\begin{problem}{BSD 的基本预测}
\easy
陈述 BSD 猜想在解析秩层面的内容，并把它应用到 $E_{11}$。
\end{problem}
\begin{solution}
BSD 的基本陈述是
\[
\ord_{s=1}L(E,s)=\operatorname{rank}E(\Q).
\]
左边叫解析秩，右边是 Mordell--Weil 群的代数秩。

对 $E_{11}$，前面章节已经知道
\[
E_{11}(\Q)\cong \Z/5\Z,
\]
因此代数秩是 $0$。BSD 于是预测
\[
\ord_{s=1}L(E_{11},s)=0,
\]
也就是说 $L(E_{11},1)\ne 0$。
\end{solution}

\begin{problem}{Weil 定理的验证}
\easy
已知 level $11$ 的唯一归一化新形式满足
\[
f_{11}(q)=q-2q^2-q^3+2q^4+q^5+2q^6-2q^7+\cdots.
\]
说明它与 $E_{11}$ 的前几个系数一致，从而给出 Weil 定理在本例中的最初验证。
\end{problem}
\begin{solution}
从点数得到
\[
a_2=-2,\quad a_3=-1,\quad a_5=1,\quad a_7=-2,
\]
又由递推算得
\[
a_4=2,\quad a_6=2.
\]
这些恰好与新形式 $f_{11}$ 的 Fourier 系数一致。因此至少在前几项上，
\[
L(E_{11},s)=\sum_{n\ge1}a_nn^{-s}
\]
与
\[
L(f_{11},s)=\sum_{n\ge1}a_nn^{-s}
\]
完全吻合。Weil 定理断言这种一致性并非偶然，而是对全部系数都成立。
\end{solution}

\section{进阶题 \medium}

\begin{problem}{局部 zeta 函数}
\medium
设 $p\nmid N$ 是好素数。证明椭圆曲线的局部 zeta 函数可写成
\[
Z(E/\mathbb{F}_p,T)=\frac{1-a_pT+pT^2}{(1-T)(1-pT)}.
\]
并把它应用到 $E_{11}$ 在 $p=5$ 的情形。
\end{problem}
\begin{solution}
对好约化椭圆曲线，Frobenius 在第一同调上的特征值为 $\alpha,\beta$，满足
\[
\alpha+\beta=a_p,
\qquad
\alpha\beta=p.
\]
Weil zeta 函数的一般形式是
\[
Z(E/\mathbb{F}_p,T)=\frac{(1-\alpha T)(1-\beta T)}{(1-T)(1-pT)}
=\frac{1-(\alpha+\beta)T+\alpha\beta T^2}{(1-T)(1-pT)}.
\]
因此得到所求公式。

对 $E_{11}$ 在 $p=5$，有 $a_5=1$，所以
\[
Z(E_{11}/\mathbb{F}_5,T)=\frac{1-T+5T^2}{(1-T)(1-5T)}.
\]
\end{solution}

\begin{problem}{素数幂系数递推}
\medium
对好素数 $p\nmid 11$，从 Euler 因子
\[
\frac{1}{1-a_pX+pX^2}=\sum_{r\ge0}a_{p^r}X^r
\]
推出递推公式
\[
a_{p^{r+1}}=a_pa_{p^r}-pa_{p^{r-1}}.
\]
并计算 $a_{25}$。
\end{problem}
\begin{solution}
把两边相乘：
\[
(1-a_pX+pX^2)\sum_{r\ge0}a_{p^r}X^r=1.
\]
比较 $X^{r+1}$ 的系数，就得到
\[
a_{p^{r+1}}-a_pa_{p^r}+pa_{p^{r-1}}=0,
\]
即
\[
a_{p^{r+1}}=a_pa_{p^r}-pa_{p^{r-1}}.
\]
对 $p=5$、$r=1$，有
\[
a_{25}=a_5^2-5=1-5=-4.
\]
\end{solution}

\begin{problem}{二次扩张点数的再计算}
\medium
不用 Frobenius 根，而只用上一题的递推公式，重新计算 $\#E_{11}(\mathbb{F}_{25})$。
\end{problem}
\begin{solution}
对好素数 $p$，Dirichlet 系数满足
\[
a_{p^2}=a_p^2-p.
\]
另一方面，Frobenius 在 $\mathbb{F}_{p^2}$ 上的迹是
\[
\alpha^2+\beta^2=a_p^2-2p=a_{p^2}-p.
\]
因此
\[
\#E(\mathbb{F}_{p^2})=p^2+1-(a_{p^2}-p).
\]
对 $E_{11}$、$p=5$，由上一题 $a_{25}=-4$，所以迹为
\[
a_{25}-5=-9.
\]
于是
\[
\#E_{11}(\mathbb{F}_{25})=25+1-(-9)=35.
\]
与基础题中的结果一致。
\end{solution}

\begin{problem}{从 Eichler 到 L 函数相等}
\medium
解释为什么 Eichler--Shimura 关系会导致
\[
L(f,s)=L(A_f,s).
\]
当 $f$ 的系数都在 $\mathbb{Q}$ 中且 $\dim A_f=1$ 时，这又为什么等同于某条椭圆曲线的 $L$-函数。
\end{problem}
\begin{solution}
Eichler--Shimura 关系把 Hecke 算子 $T_p$ 与 Frobenius 作用联系起来，因此对每个好素数 $p\nmid N$，$T_p$ 的本征值 $a_p(f)$ 就等于相应 Abelian variety $A_f$ 上 Frobenius 的迹。于是两边的局部 Euler 因子相同：
\[
1-a_p(f)p^{-s}+p^{1-2s}.
\]
坏素数处也有一致的局部描述，所以乘起来得到
\[
L(f,s)=L(A_f,s).
\]
若 $f$ 的 Fourier 系数都落在 $\mathbb{Q}$ 中，而且 $\dim A_f=1$，那么 $A_f$ 本身就是一条定义在 $\mathbb{Q}$ 上的椭圆曲线 $E$。于是
\[
L(f,s)=L(A_f,s)=L(E,s).
\]
这就是 Weil 定理在维数 $1$ 情形下的内容。
\end{solution}

\begin{problem}{函数方程的符号}
\medium
设完成 $L$-函数满足
\[
\Lambda(E,s)=\varepsilon\Lambda(E,2-s),\qquad \varepsilon=\pm1.
\]
说明当 $\varepsilon=-1$ 时一定有 $L(E,1)=0$；当 $\varepsilon=+1$ 时，解析秩必须是偶数。
\end{problem}
\begin{solution}
把 $s=1$ 代入函数方程，得到
\[
\Lambda(E,1)=\varepsilon\Lambda(E,1).
\]
若 $\varepsilon=-1$，则
\[
\Lambda(E,1)=-\Lambda(E,1),
\]
所以 $\Lambda(E,1)=0$，也即 $L(E,1)=0$。

若 $\varepsilon=+1$，则 $\Lambda(E,s)$ 关于 $s=1$ 是偶对称的，因此在 $s=1$ 的零点阶数必须是偶数。换言之，解析秩一定是偶数。
\end{solution}

\begin{problem}{BSD 的精确公式}
\medium
写出 BSD 猜想在秩零情形的精确公式，并解释其中 $\Sha(E)$、$c_p$、$\Omega_E$ 和挠子群各自扮演的角色。
\end{problem}
\begin{solution}
当 $\operatorname{rank}E(\Q)=0$ 时，BSD 的精确公式写成
\[
\frac{L(E,1)}{\Omega_E}=\frac{\#\Sha(E)\prod_p c_p}{\#E(\Q)_{\mathrm{tors}}^2}.
\]
其中：
\begin{itemize}
  \item $\Omega_E$ 是实周期，衡量椭圆曲线的解析大小；
  \item $\Sha(E)$ 是 Tate--Shafarevich 群，反映局部到整体的障碍；
  \item $c_p$ 是 Tamagawa 数，记录坏约化处 N\'eron 模型分量群的贡献；
  \item $E(\Q)_{\mathrm{tors}}$ 的大小出现在分母，说明有理挠点越多，中心值比值往往越小。
\end{itemize}
这是 BSD 在秩零时最直接可检验的形式。
\end{solution}

\begin{problem}{导子十一曲线的 BSD 预测}
\medium
对 $E_{11}$ 利用已知数据
\[
E_{11}(\Q)_{\mathrm{tors}}\cong \Z/5\Z,
\qquad c_{11}=5,
\qquad \operatorname{rank}E_{11}(\Q)=0,
\]
把 BSD 精确公式化简成只含 $\Sha(E_{11})$ 的形式。
\end{problem}
\begin{solution}
对 $E_{11}$，所有 Tamagawa 数中只有 $c_{11}$ 非平凡，所以
\[
\prod_p c_p=5.
\]
又因为挠子群阶为 $5$，BSD 秩零公式化为
\[
\frac{L(E_{11},1)}{\Omega_{E_{11}}}
=\frac{\#\Sha(E_{11})\cdot 5}{5^2}
=\frac{\#\Sha(E_{11})}{5}.
\]
因此只要数值上知道左边，就能直接预测 $\Sha(E_{11})$ 的大小。特别地，若数值比值接近 $1/5$，就强烈暗示
\[
\#\Sha(E_{11})=1.
\]
\end{solution}

\begin{problem}{中心值的快速近似}
\medium
利用权 $2$ 新形式的 Mellin 变换公式
\[
L(E_{11},1)=2\sum_{n\ge1}\frac{a_n}{n}e^{-2\pi n/\sqrt{11}},
\]
用前六项系数估计 $L(E_{11},1)$。
\end{problem}
\begin{solution}
记
\[
q=e^{-2\pi/\sqrt{11}}\approx 0.1504007847.
\]
用
\[
a_1=1,a_2=-2,a_3=-1,a_4=2,a_5=1,a_6=2
\]
得到
\[
L(E_{11},1)\approx 2\left(q-\frac{2q^2}{2}-\frac{q^3}{3}+\frac{2q^4}{4}+\frac{q^5}{5}+\frac{2q^6}{6}\right).
\]
数值计算为
\[
2\left(0.1504008-0.0226204-0.0011340+0.0002558+0.0000154+0.0000039\right)
\approx 0.25384.
\]
由于 $q\approx 0.15$ 很小，后面的项衰减极快，所以这个近似已经相当精确，支持 $L(E_{11},1)\ne 0$。
\end{solution}

\begin{problem}{部分 Euler 乘积}
\medium
对 $E_{11}$ 定义
\[
P(X)=\prod_{p\le X,\ p\ne11}\frac{p}{\#E_{11}(\mathbb{F}_p)}.
\]
利用第 7 章的数据计算 $P(5)$、$P(13)$、$P(29)$，并解释这与秩零情形为何相符。
\end{problem}
\begin{solution}
第 7 章已经算得
\[
P(5)=\frac25\cdot\frac35\cdot\frac55=0.24,
\]
\[
P(13)=0.24\cdot\frac7{10}\cdot\frac{13}{10}=0.2184,
\]
\[
P(29)=0.2184\cdot\frac{17}{20}\cdot\frac{19}{20}\cdot\frac{23}{25}\cdot\frac{29}{30}
\approx 0.15684.
\]
这些值没有出现持续上升，更像稳定在常数量级。对秩零曲线，Birch--Swinnerton-Dyer 的原始数值实验正预期这种有界行为，因此它与 $E_{11}$ 的秩零结论相吻合。
\end{solution}

\section{挑战题 \hard}

\begin{problem}{中心值与解析秩}
\hard
综合函数方程和上一题的数值估计，说明为什么 $E_{11}$ 的解析秩应当是零。
\end{problem}
\begin{solution}
由完成 $L$-函数的函数方程，$E_{11}$ 的解析秩是中心点 $s=1$ 处的零点阶数。前一题给出的快速级数估计说明
\[
L(E_{11},1)\approx 0.25384\ne0.
\]
因此 $s=1$ 不是零点，零点阶数只能是
\[
\ord_{s=1}L(E_{11},s)=0.
\]
这就表明解析秩为零。再与第 6 章中 $E_{11}(\Q)$ 的代数秩为零相比较，可见 BSD 在这个例子里得到完全一致的数值支持。
\end{solution}

\begin{problem}{与导子三十七的对比}
\hard
设 $E_{37}$ 是一条导体为 $37$ 的模椭圆曲线，且其函数方程符号为 $-1$。证明 $L(E_{37},1)=0$，并说明 BSD 因而预测什么样的秩。
\end{problem}
\begin{solution}
函数方程写成
\[
\Lambda(E_{37},s)=-\Lambda(E_{37},2-s).
\]
令 $s=1$，得到
\[
\Lambda(E_{37},1)=-\Lambda(E_{37},1),
\]
故
\[
\Lambda(E_{37},1)=0.
\]
所以 $L(E_{37},1)=0$。这说明解析秩至少为 $1$；又因为符号为 $-1$ 时中心点附近是奇对称，零点阶数必须是奇数。BSD 因此预测
\[
\operatorname{rank}E_{37}(\Q)
\]
是正的奇数；在最常见的情形下就是秩 $1$。
\end{solution}

\begin{problem}{更高次扩张上的点数}
\hard
设 $S_n=\alpha^n+\beta^n$，其中 $\alpha+\beta=a_p$、$\alpha\beta=p$。证明
\[
S_n=a_pS_{n-1}-pS_{n-2}
\]
并用它计算 $\#E_{11}(\mathbb{F}_8)$。
\end{problem}
\begin{solution}
因为 $\alpha,\beta$ 都满足二次方程
\[
X^2-a_pX+p=0,
\]
所以对任意 $n\ge2$，
\[
\alpha^n=a_p\alpha^{n-1}-p\alpha^{n-2},
\qquad
\beta^n=a_p\beta^{n-1}-p\beta^{n-2}.
\]
相加即得
\[
S_n=a_pS_{n-1}-pS_{n-2}.
\]
对 $p=2$，有 $a_2=-2$，并且
\[
S_0=2,
\qquad S_1=-2,
\qquad S_2=a_2^2-2p=0.
\]
于是
\[
S_3=a_2S_2-2S_1=(-2)\cdot 0-2(-2)=4.
\]
因此
\[
\#E_{11}(\mathbb{F}_8)=2^3+1-S_3=8+1-4=5.
\]
\end{solution}

\begin{problem}{Euler 乘积与系数一致性}
\hard
只使用 $p=2,3$ 两个局部因子，展开
\[
\frac{1}{1-a_22^{-s}+2^{1-2s}}\cdot \frac{1}{1-a_33^{-s}+3^{1-2s}}
\]
到包含 $2^{-2s}$、$3^{-2s}$ 和 $6^{-s}$ 的精度，检查得到的系数与 $a_2,a_3,a_4,a_6,a_9$ 是否一致。
\end{problem}
\begin{solution}
先看 $p=2$。因为 $a_2=-2$，局部因子是
\[
\frac{1}{1+2\cdot 2^{-s}+2\cdot 2^{-2s}}=1+a_2 2^{-s}+a_4 2^{-2s}+\cdots,
\]
其中
\[
a_4=a_2^2-2=2.
\]
同理对 $p=3$，由 $a_3=-1$ 得
\[
\frac{1}{1+3^{-s}+3\cdot 3^{-2s}}=1+a_3 3^{-s}+a_9 3^{-2s}+\cdots,
\]
且
\[
a_9=a_3^2-3=-2.
\]
把两者相乘，只保留到要求的精度：
\[
1+a_2 2^{-s}+a_3 3^{-s}+a_4 2^{-2s}+a_6 6^{-s}+a_9 3^{-2s}+\cdots.
\]
其中交叉项系数是
\[
a_6=a_2a_3=(-2)(-1)=2.
\]
所以展开系数正好是
\[
a_2=-2,\quad a_3=-1,\quad a_4=2,\quad a_6=2,\quad a_9=-2,
\]
与前面由 Hecke 递推得到的数值一致。
\end{solution}

\begin{problem}{坏素数的局部因子}
\hard
说明 $E_{11}$ 在 $11$ 处的分裂乘法约化为什么对应局部因子
\[
(1-11^{-s})^{-1},
\]
并把它与好素数处的二次局部因子作比较。
\end{problem}
\begin{solution}
对好素数 $p\nmid N$，Frobenius 在第一同调上有两个非平凡特征值，因此局部因子是二次的：
\[
(1-a_pp^{-s}+p^{1-2s})^{-1}.
\]

而在乘法坏约化处，几何对象退化成一个结点曲线；同调上只剩一个有效的一次局部参数，因此局部因子退化成一次项。对分裂乘法约化，符号取正，于是
\[
L_{11}(E_{11},s)=(1-11^{-s})^{-1}.
\]
若是非分裂乘法约化，则会变成 $(1+11^{-s})^{-1}$。所以坏素数处局部因子次数下降，正反映了曲线几何上的退化。
\end{solution}

\begin{problem}{BSD 比值与沙群预测}
\hard
已知第 7 章给出 $c_{11}=5$、$E_{11}(\Q)_{\mathrm{tors}}\cong \Z/5\Z$。若数值上观察到
\[
\frac{L(E_{11},1)}{\Omega_{E_{11}}}\approx 0.2,
\]
说明 BSD 会如何预测 $\Sha(E_{11})$。
\end{problem}
\begin{solution}
在进阶题中已经化简得到
\[
\frac{L(E_{11},1)}{\Omega_{E_{11}}}=\frac{\#\Sha(E_{11})}{5}.
\]
若数值上左边接近 $0.2=1/5$，那么最自然的预测是
\[
\frac{\#\Sha(E_{11})}{5}=\frac15.
\]
于是
\[
\#\Sha(E_{11})=1.
\]
也就是说，BSD 在这个例子里预言 Tate--Shafarevich 群是平凡的。
\end{solution}

\begin{problem}{更长的数值证据}
\hard
继续计算部分乘积
\[
P(X)=\prod_{p\le X,\ p\ne11}\frac{p}{\#E_{11}(\mathbb{F}_p)}
\]
到 $X=97$，已知数值约为 $P(97)\approx 0.15414$。结合 $P(5),P(13),P(29)$，解释这个现象支持怎样的 BSD 图景。
\end{problem}
\begin{solution}
我们已经有
\[
P(5)=0.24,
\qquad
P(13)=0.2184,
\qquad
P(29)\approx 0.15684,
\qquad
P(97)\approx 0.15414.
\]
这些数值并没有像对数那样持续增长，而是快速落入一个稳定区间，随后只做温和波动。这正是原始 BSD 数值实验在秩零情形下期待看到的行为：相关乘积应近似常数，而不是发散。于是它与
\[
\operatorname{rank}E_{11}(\Q)=0
\]
和
\[
L(E_{11},1)\ne0
\]
的图景彼此一致。
\end{solution}

\begin{problem}{从 Eichler到 Weil}
\hard
用一句完整的逻辑链说明：为什么对 $E_{11}$ 而言，Eichler--Shimura 关系最终会导出
\[
L(E_{11},s)=L(f_{11},s).
\]
\end{problem}
\begin{solution}
逻辑链如下：
\begin{itemize}
  \item[(i)] 第 6 章说明 $X_0(11)$ 的 Jacobian 在维数 $1$ 情形下就是一条椭圆曲线，可取为 $E_{11}$；
  \item[(ii)] 第 8 章的 Eichler--Shimura 关系把 $T_p$ 与 Frobenius 的迹联系起来，因此新形式 $f_{11}$ 的 Hecke 本征值 $a_p(f_{11})$ 恰是 $E_{11}$ 上 Frobenius 的迹；
  \item[(iii)] 所以对每个好素数，$f_{11}$ 和 $E_{11}$ 的局部 Euler 因子相同；坏素数处也由约化类型给出一致的局部因子；
  \item[(iv)] 把所有局部因子相乘，就得到全局恒等式
  \[
  L(E_{11},s)=L(f_{11},s).
  \]
\end{itemize}
这正是 Weil 定理在导体 $11$ 例子中的具体体现。
\end{solution}

\begin{problem}{第八码内容的综合理解}
\hard
综合本章内容，说明为什么 $E_{11}$ 是研究 BSD 猜想的理想入门例子。
\end{problem}
\begin{solution}
$E_{11}$ 之所以是理想入门例子，有四个原因。
\begin{itemize}
  \item[(i)] 它来自最简单的正亏格模曲线 $X_0(11)$，几何背景清晰；
  \item[(ii)] 它的 Fourier 系数可以由小素数点数直接算出，Hecke 本征值与 Frobenius 迹的关系非常具体；
  \item[(iii)] 完成 $L$-函数有明确的函数方程，并能用快速收敛级数直接估计中心值；
  \item[(iv)] 它的代数数据也很简单：导体 $11$、Tamagawa 数 $5$、有理挠子群 $\Z/5\Z$、秩为 $0$，因此 BSD 精确公式化简成
  \[
  \frac{L(E_{11},1)}{\Omega_{E_{11}}}=\frac{\#\Sha(E_{11})}{5}.
  \]
\end{itemize}
于是从模曲线、Hecke 算子、Eichler--Shimura、函数方程到 BSD 数值检验，都能在一条具体曲线上完整看到。
\end{solution}
